\begin{figure}
\begin{lstlisting}
Lemma poly_multiply
  (n p r h r0 r1 h0 h1 h2 s1 d0 d1 d2 hh : Z)
  (H: p > 0 /\ r1 >= 0 /\ n > 0 /\ 4 * (n * n) = p + 5 /\ r = r1 * n + r0 /\
            h = h2 * (n * n) + h1 * n + h0 /\ s1 = r1 + (r1 / 4) /\ r1 mod 4 = 0 /\
            d0 = h0 * r0 + h1 * s1 /\ d1 = h0 * r1 + h1 * r0 + h2 * s1 /\
            d2 = h2 * r0 /\ hh = d2 * (n * n) + d1 * n + d0)
  : ((h * r) mod p = hh mod p).
Proof.
  repeat match goal with H : ?A /\ ?B |- _ => destruct H end.
  subst.
  match goal with |- ?A mod p = ?B mod p => change (eqm p A B) end.
  match goal with K : r1 mod 4 = 0 |- _ =>
   apply Z_div_exact_full_2 in K; try omega;
   revert K;
   generalize (r1 / 4);
   intro u;
   intros; subst
  end.
  pose (b := (h2 * n + h1) * u).
  generalize (Zopp_eqm p).
  generalize (Zplus_eqm p).
  generalize (Zminus_eqm p).
  generalize (Zmult_eqm p).
  generalize (eqm_setoid p).
  set (e := eqm p).
  intros.
  match goal with |- e _ ?R =>
    generalize (modulo_addition_lemma R p b)
  end.
  fold e. intro K. setoid_rewrite <- K. clear K.
  assert (p = 4 * (n * n) - 5) as J by omega.
  replace (b * p) with (b * (4 * (n * n) - 5)) by congruence.
  unfold b.
  apply eq_implies_eqm.
  ring.
Qed.
\end{lstlisting}
\caption{Coq proof of \ls$poly_multiply$}
\label{fig:coqpolymultiply}
\end{figure}

\section{Metaprogramming verified parsers and serializers}
\label{sec:parsers-detail}

We consider parser and serializer specifications as ghost code that
``operates'' on finite sequences of bytes of unbounded lengths.
For a given type $t$, a serializer for $t$ marshals any data of type
$t$ into a finite sequence of bytes; and a parser for $t$ reads from a
sequence of bytes, checks whether it corresponds to a valid data of
type $t$, and if so, returns that parsed data and the number of bytes
consumed.

The serializer specification always succeeds; the parser specification
succeeds if and only if the input data is valid with respect to the
parser. On those specifications, we aim to prove that the parser and
the serializer are partial inverses of each other. We encode these
requirements as refinements on the types of parser and serializer
specifications:
%
\begin{lstlisting}
type byte = FStar.UInt8.t
type parser t = (input: seq byte) -> GTot (option (t * (x: nat { x < length input } )))
let serialize_then_parse_eq_id #t (p: parser t) (s: (t -> GTot (seq byte))) =
  (forall (x: t) . let y = s x in p y == Some (x, length y))
let parse_then_serialize_eq_id #t (p: parser t) (s: (t -> GTot (seq byte))) =
  (forall (x: seq byte) . match p x with | Some (y, len) -> s y == slice x 0 len | _ -> True) })
type serializer #t (p: parser t) = (s: (t -> GTot (seq byte))
  { serialize_then_parse_eq_id p s /\ parse_then_serialize_eq_id p s })
\end{lstlisting}

Whereas those specifications are ghost code, we would like to generate
implementations that can be extracted to C. To this end, we generate
\emph{stateful} implementations written in the \lowstar set of \fstar~\cite{lowstar},
operating on buffers, which are \lowstar mutable data structures
representing C arrays. There, instead of a sequence of bytes, the
parser implementation is given an input buffer and its length; and the
serializer implementation is given an output buffer (along with its
length) onto which it is to serialize the data. This means that, given
a piece of data to serialize and a destination buffer, a serializer
implementation can succeed only if the serialized data fits into the
destination buffer; if so, then the serializer will return the number
of bytes written.

We bake the correctness of the implementations with respect to their
specifications at the level of their types:

\begin{lstlisting}
type buffer8 = LowStar.Buffer.buffer FStar.UInt8.t
type u32 = FStar.UInt32.t

let parser32_postcond #t (p: parser t) (input: buffer8) (l: u32 {l == len input})
  (h: mem) (res: option (t * u32)) (h': mem) =
  modifies loc_none h h' /\ (* memory safety *)
  match p (as_seq h input), res with (* functional correctness *)
  | Some (x, ln), Some (x', ln') -> x' == x /\ FStar.UInt32.v ln' == ln
  | None, None -> True   | _ -> False

type parser32 #t (p: parser t) = (input: buffer8) -> (l: u32 { l == len input } )
                                              -> Stack (option (t * u32))
  (requires (fun h -> live h input)) (ensures (parser32_postcond p input))

let serializer32_postcond #t (p: parser t) (s: serializer p) (output: buffer8)
  (l: u32 {l == len output}) (x: t) (h: mem) (res: option u32) (h': mem) =
  live h output /\ live h' output /\ modifies (loc_buffer b) h h' /\ (* memory safety *)
  let ln = length (s x) in
  match res with (* functional correctness *)
  | Some ln' ->
    FStar.UInt32.v ln' == ln /\ ln <= FStar.UInt32.v l /\
    let b' = sub b 0ul ln' in modifies (loc_buffer b') h h' /\ as_seq h' b' == s x
  | None -> ln > FStar.UInt32.v l

type serializer32 #t (#p: parser t) (s: serializer p) =
  (output: buffer8) -> (l: u32 { l == len output } ) -> (x: t) -> Stack (option u32)
  (requires (fun h -> live h output)) (ensures (serializer32_postcond s output l x))
\end{lstlisting}

\tahina{TODO: gen\_parser can be implemented using typeclasses, triggering gen\_enum\_parser.}

\paragraph{\bf Generating the specification from a \fstar datatype}

Instead of mandating the user to write the parser and serializer
specification themselves, we write a \metafstar
tactic, \verb+gen_specs+, to generate both
the parser and the serializer specifications directly from the \fstar
type $t$ given by the user. These tactics operate by syntax inspection
on $t$ itself, and generates the parser and serializer specifications
using our combinators.

Contrary to Coq's Ltac, \metafstar tactics can also inspect
definitions of inductive types. From there, we made \verb+gen_specs+
also generate a parser and serializer specification out of an
enumeration type defined as a \fstar inductive type whose constructors
have no arguments: for such a type with $n \leq 256$ constructors, the
corresponding parser will associate a unique 8-bit integer from $0$ to
$n-1$ to each constructor. For instance, if $t$ is defined as a \fstar
enumeration type:

As before, by virtue of the correctness of serializers being baked in
their type, the correctness of the serializer generated
by \verb+gen_specs+ does not appear as such as a verification
condition when \fstar is type-checking the generated term. In the
example above, the only proof obligations generated are those to
correctly type-check the rewriting functions themselves passed
to \verb+serialize_synth+ (e.g. enum constructors are rewritten to
integers less than 4), and the precondition of \verb+serialize_synth+
(i.e. the fact that the two rewriting functions are inverse of each
other.) These can be discharged by the SMT solver, but they can also
be discharged automatically by our tactics directly by case analysis
on the inductive type, or by bounded integer enumeration.

\paragraph{\bf Binders}

For instance, we specified \verb+seq_p+, a combinator for dependent
parsing:\protect{\footnote{This parsing combinator is based on yet
unpublished work by Tej Chajed. Building serializers for parsers
derived from \ls$seq_p$ is explicitly out of the scope of this
paper.}} \tahina{I mentioned Tej in the footnote. This part of the
footnote should be removed if Tej becomes a coauthor of the paper.}

\begin{lstlisting}
let seq_p #t1 (p1: parser t1) #t2 (p2: (t1 -> Tot (parser t2))) : Tot (parser t2) =
  fun input -> match p1 input with | None -> None | Some (x1, cons1) ->
    let input2 = slice input cons1 (length input) in
    match p2 x1 input2 with
    | None -> None | Some (x2, cons2) -> Some (x2, cons1 + cons2)
\end{lstlisting}

Then, we implemented it in \lowstar (\verb+seq_pi+), and
manually proved it correct. From there, the user can write a simple
parser specification for a parser that parses an integer encoded
in one or two bytes depending on its value (where \verb+parse_ret x+
is a parser that returns \verb+x+ without reading its byte input:)

\begin{lstlisting}
let example : parser_spec FStar.Int16.t =
  seq_p parse_u8
           (fun lo -> if lo <^ 128uy then parse_ret (cast_u8_to_i16 lo)
                   else parse_synth parse_u8 (fun hi -> cast_u8_to_i16 (lo %^ 128uy)
                                                 +^ cast_u8_to_i16 hi))
\end{lstlisting}
%
Then, if the user writes:
%
\begin{lstlisting}
let example_impl : parser_impl example = _ by gen_parser_impl
\end{lstlisting}
%
then \verb+gen_parser32+ inspects the shape of the \emph{goal},
which is \verb+parser32 example+, and so generates the following \lowstar implementation:
%
\begin{lstlisting}
seq_pi parse32_u8 (fun lo ->
  parse32_ifthenelse (lo <^ 128uy) (fun _ -> parse32_ret (cast_u8_to_i16 lo)) (fun _ -> parse32_synth parse32_u8 (fun hi -> cast_u8_to_i16 (lo %^ 128uy) +^ cast_u8_to_i16 hi)))
\end{lstlisting}

from where KreMLin \cite{lowstar} then inlines the corresponding implementation combinators to produce C code.

First, \verb+gen_parser32+ inspects the shape of the goal to determine
the type of the data to be parsed and the parser specification; then,
it calls a recursive tactic \verb+gen_parser32'+ that inspects the
shape of the parser specification and actually builds the
implementation.

\begin{lstlisting}
module T = FStar.Tactics
let rec gen_parser32' (env: T.env) (t: T.term) (p: T.term) : T.Tac T.term =
  let (hd, tl) = app_head_tail p in
  if hd `T.term_eq` (`(parse_ret)) then T.mk_app (`(parse32_ret)) tl else
  if hd `T.term_eq` (`(parse_u8)) then (`(parse32_u8)) else
  if hd `T.term_eq` (`(seq_p)) then match tl with
  | [(t, _); (p, _); (t', _); (p', _)] ->
    begin match T.inspect p' with
    | T.Tv_Abs bx body ->
      let p32 = gen_parser32' env k t p in
      let env' = T.push_binder env bx in
      let body' = gen_parser32' env' k' t' body in
      let p32' = T.pack (T.Tv_Abs bx body') in
      (`(seq_pi #(`#t) #(`#p) (`#p32) (`#t') (`#p') (`#p32')))
    | _ -> ...
    end
  | _ -> ...
  else
  ...

let gen_parser32 () : T.Tac unit =
  let (hd, tl) = app_head_tail (T.cur_goal ()) in
  if hd `T.term_eq` (`parser32)
  then match tl with
  | [(t, _); (p, _)] ->
    let env = T.cur_env () in
    let p32 = gen_parser32' env t p in
    T.exact_guard p32;
  | _ -> ...
\end{lstlisting}

When \verb+gen_parser32+ (resp. \verb+gen_serializer32+) builds the
parser (resp. serializer) implementation, it may generate verification
conditions due to specific preconditions required by certain
combinators (such as the precondition of the \verb+serialize_synth+
rewriting serializer combinator: the rewriting functions being inverse
of each other). Contrary to Coq, there are no proof objects, so those
preconditions need to be proven again even if they had been proven
earlier by the user at the specification level. Nevertheless, those
preconditions can be automatically solved by our tactics, or they can
directly be sent to the SMT solver.

However, by virtue of the correctness of the implementation
combinators being baked in their type, typechecking the resulting term
generated by \verb+gen_parser32+ triggers no verification condition
other than those preconditions required by the specific combinators,
and unification constraints, the latter solved by reflexivity. In
particular, the memory safety and correctness of the implementations
generated by \verb+gen_parser32+ do not appear as such as verification
conditions when \fstar is type-checking the terms generated
by \verb+gen_parser32+. \tahina{Is it already important to mention
here the kinds of proof obligations generated by the tactics, and that
they can be solved either by SMT or tactics? Or should we move to
section 5?}

\subsection{Verified meta-programming of program transformations}

\ls$compile_bind$ expects, in \ls$bind f1 f2$, that \ls$f2$ be an
abstraction, and so needs to manipulate binders and recursively
compile \ls$f1$ and the body of \ls$f2$. This will, in the latter
case, make some bound variables appear free in terms being compiled.

If the head term $t_0$ is a free variable not corresponding to any
combinator, then \ls$compile_fvar$ unfolds it, compiles the unfolded
term, and inserts a coercion.\footnote{We defined an explicit coercion
combinator, \ls$coerce_sz$, such that \ls$coerce_sz f1 f1_sz f2$ is
well-typed and is a \ls$m_sz f2$ as soon as \ls$f1 () == f2 ()$ as $m$
specifications.} In most cases, a free variable to be unfolded will
correspond to a function definition, so that unfolding will yield a
$\beta$-redex, which \ls$compile_fvar$ reduces by calling a
normalization tactic that we provide as part of \metafstar. This needs
an \emph{environment}
\ls{(e: env)} tracking the bound variables encountered by nested calls
to \ls$compile_bind$.

We show below the implementation of \ls$compile_bind$, which is
representative of most features of \metafstar that we use for
our \ls$compile$ tactic.

\begin{lstlisting}
 (* extract the head symbol and the arguments of an application *)
let rec app_head_rev_tail (t: term) : Tac (term * list argv) = 
  match inspect t with | Tv_App u v -> let (x, l) = app_head_rev_tail u in (x, v :: l) | _ -> (t, [])
let app_head_tail (t: term) : Tac (term * list argv) = let (x, l) = app_head_rev_tail t in (x, rev l)

let compile_bind (e: env) (ty: term) (t: term) (compile: env -> term -> term -> Tac term) : Tac term =
  let (t_0, ar) = app_head_tail t in
  guard (t_0 = quote bind); (* fail if false *)
  match ar with
  | [ ($\tau_1$, _); ($\tau_2$, _); (f1, _); (f2, _); ] ->
    begin match inspect f2 with
    | Tv_Abs v f2_body ->
      let f1' = compile e $\tau_1$ f1 in
      let e' = push_binder e v in
      let f2_body' = compile e' $\tau_2$ f2_body in
      let f2' = pack (Tv_Abs v f2_body') in
      (`(bind_sz #(`#tau_1) #(`#tau_2) (`#f1) (`#f1') (`#f2) (`#f2')))
    | _ -> fail "compile_bind: Not an abstraction"
    end
  | _ -> fail "compile_bind: 4 arguments expected"
\end{lstlisting}

\paragraph{\bf Putting everything together}
Finally, the following wrapper combinator $\Phi$ puts everything
together, taking a specification $f$ and its size and buffer
implementations, and yielding a stateful function, fitting in
the \lowstar subset of \fstar amenable to extraction to C, that first
computes the expected size of the output, then allocates one single
buffer using the C \ls$malloc$, then uses it to write the output of
$f$:
\begin{lstlisting}
let $\Phi$ $\tau$ (f: m $\tau$) (f_sz: m_sz f) (f_st: m_st f) : ST (option ($\tau$ * buffer U8.t)) (requires (fun _ -> True))
  (ensures (fun h res h' -> let (r, log) = f () in
    if length log > $2^{31} - 1$ then res == None else
    Some? res /\ (let (Some (r', b)) = res in r == b /\ fresh b h h' /\  as_seq h b == log)
  )) = match f_sz with | None -> None | Some (_, sz') ->
       let b = rcreate_mm root 42uy sz in (* then allocate a fresh buffer *)
       let (r, _) = f_st b in (* then write the output into the buffer *)
       Some (r, b)
\end{lstlisting}

%%% Tahina: this paragraph removed because it should rather go to a
%%% parser-specific paper.
%%% CH: Added this back, Nada is on the ESOP %%% PC!
\paragraph{\bf Related Work}
Our approach of generating verified low-level formatters is related to
Amin and Rompf's~\cite{Amin:2017:LAW:3009837.3009867} work on
LMS-Verify. They use metaprogramming facilities in Scala to generate C
code together with proof annotations to be checked by
Frama-C~\citep{CuoqKKPSY12}, an SMT-based C verifier. In contrast, we
generate correct-by-construction imperative \lowstar code, which is
verified by \fstar{} before being translated to C by the KreMLin
compiler.  LMS-Verify has also been used to generate efficient and
safe HTTP parsers, a larger-scale effort than the network message
formatters than we have currently done.

%%%%% Appendix for Section 4

\newpage

\section{Proof of \autoref{thm:split}}
\label{app:correctness}
This proof is based on the formal definition of \emf, a recent formalization
of an \fstar{} subset~\cite{dm4free}.
%x
We use the same notation and rule names from there, but the proof is
nevertheless quite direct with just minimum familiarity with \emf.
%
We use $E$ to represent \fstar contexts:

\[
\begin{array}{rclcl}
    E &=&    \cdot  & & \\
      & \mid & t E &                 \mid & E t \\
      & \mid & \lambda(x:t). E &     \mid & \lambda(x:E). t \\
      & \mid & (x:t) \rightarrow E & \mid & (x:E) \rightarrow t \\
      & \mid & x:t\{E\} &            \mid & x:E\{t\} \\
      & ~ & \ldots
\end{array}
\]

Contexts present a set of bound variables to their hole. We denote
those variables by $\gamma(E)$:

\[
\begin{array}{lcl}
    \gamma(\cdot) &=& \emptyset \\
    \gamma(\lambda(x:t). E) &=& \{x\} \cup \gamma(E) \\
    \gamma(\lambda(x:E). t) &=& \gamma(E) \\
    \ldots
\end{array}
\]
    % γ(∘) = {}
    % γ(λ(x:t). E) = (x:t) ++ γ(E)
    % γ(λ(x:E). t) = γ(E)
    % ...

We say a context $E$ has type $t_1 \Rightarrow t_2$ for $\Gamma$ (noted
by $\Gamma \vdash E : t_1 \Rightarrow t_2$) when for all $e$ such that
$\Gamma, \gamma(E) \vdash e : t_1$ we also have $\Gamma \vdash E[e] : t_2$.

\subsection{Soundness of splitting the proof obligation}

The proof follows mainly from the following lemma:

\begin{lemma}
    \[
    \infer{\Gamma \vDash E[t_1] = E[t_2]}
          {\Gamma, \gamma(E) \vDash t_1 = t_2}
    \]
\end{lemma}
\begin{proof}
    The proof follows from applying functional extensionality and using
    a carefully crafted abstraction.

    Let $S = \lambda f. E[f (\gamma(E))]$ (where $f (\gamma(E))$
    represents $f$ applied to every variable in $\gamma(E)$,
    sequentially). Also, $\lambda \gamma(E). t$ represents
    abstracting $E$ over the variables in $\gamma(E)$, and
    similarly for $\forall \gamma(E). t$.

    Then,
    \[
        \infer[\mbox{Reduction}]{\Gamma \vDash E[t_1] = E[t_2]}
        {\infer[\mbox{V-EqP}]{\Gamma \vDash S (\lambda \gamma(E). t_1) = S (\lambda \gamma(E). t_2)}
        {\infer[\mbox{V-Ext}]{\Gamma \vDash (\lambda \gamma(E). t_1) = (\lambda \gamma(E). t_2)}
        {\infer[\mbox{V-$\forall$i}]{\Gamma \vDash \forall \gamma(E). t_1 = t_2}
                             {\Gamma, \gamma(E) \vDash t_1 = t_2}}}}
    \]

    Note that the particular set of contexts E does not influence
    the proof.



        % Γ, γ(E) |= t1 = t2
        % ------------------------- [V-∀i, repeatedly]
        % Γ |= forall γ(E). t1 = t2
        % ------------------------------ [V-Ext, repeatedly]
        % Γ |= (λγ(E). t1) = (λγ(E). t2)
        % --------------------------------- [V-EqP, with P = S]
        % Γ |= S (λy(E). t1) = S (γ(E). t2)
        % --------------------------------- [Red]
        % Γ |= E[t1] = E[t2]
\end{proof}

Having such lemma, the theorem is proven as follows:
\[
    \infer{\Gamma \vDash E[\phi]}
         {\Gamma \vDash E[\top]
         & \infer[\mbox{Lemma}]{\Gamma \vDash E[\phi] = E[\top]}{
         \infer[\mbox{PropExt}]{\Gamma, \gamma(E) \vDash \phi = \top}{
         \infer[\mbox{Trivial}]{\Gamma, \gamma(E) \vDash \phi \iff \top}
                               {\Gamma, \gamma(E) \vDash \phi}}}
         }
\]


\subsection{Partial completeness of splitting the proof obligation}

While the previous theorem allows to soundly split any
subformulae within a VC, we have seen that some of them
constraint the system.
%
Here we prove that for a particular set of contexts, splitting
does not make the judgments stronger, and via Theorem 1, then they are
equivalent.

We define a particular shape of ``positive'' contexts $P$.
We don't attempt to follow all of EMF*'s syntax, just the parts
we commonly see as VCs in practice.

\[
\begin{array}{rclcl}
    P &=&    \cdot  & & \\
       & \mid & \phi \land P            &\mid& P \land \phi \\
       & \mid & \forall (x:t). P \\
\end{array}
\]

In \emf, an implication $a \Rightarrow b$ is just sugar for
$\forall (\_:a). b$, so they are considered as well.

% This is sufficient to catch most VCs we encounter. Most likely, a reason
% is that disjunction/existentials can never appear in strictly-positive
% positions due to the inherent conjunctivity of WPs, but we haven't
% proved it. In any case, this property is not crucial in any way.

\begin{theorem}
    Take $\Gamma$, $E$, and $\phi$ such that:
    (1) $\phi$ is squashed,
    (2) $\Gamma \vdash E : prop \rightarrow prop$
    and (3) $\Gamma \vdash \phi : prop$.
    Then the following holds

    \[
        \infer{\Gamma, \gamma(P) \vDash \phi \qquad \Gamma \vDash P[\top]}
              {\Gamma \vDash P[\phi]}
    \]
\end{theorem}
\begin{proof}
    By induction on the structure of $P$, keeping $\Gamma$ and $\phi$
    universally quantified.
    \begin{itemize}
        \item For $P = \cdot$, trivial
        \item For $P = \psi \land P'$, our hypothesis is
              $\Gamma \vDash \psi \land P'[\phi]$.
              %
              Hence,
              $\Gamma \vDash \psi$ and
              $\Gamma \vDash P'[\phi]$.
              %
              By the IH, we get
              $\Gamma, \gamma(P') \vDash \phi$
              (note that $\gamma(P) = \gamma(P')$)
              and
              $\Gamma \vDash P'[\top]$.
              %
              By weakening and conjunction, 
              we can prove $\Gamma, \gamma(P') \vDash \psi \land \phi$
              and $\Gamma \vDash \psi \land P'[\top]$, and conclude.

        \item For $P = P' \land \psi$, the reasoning is analogous.

        \item For $P = \forall (x:t). P'$, our hypothesis is
              $\Gamma \vDash \forall (x:t). P'[\phi]$. Hence,
              we get $\Gamma, x:t \vDash P'[\phi]$ by
              eliminating the quantifier.
              %
              From the IH, we get
              $\Gamma, x:t, \gamma(P') \vDash \phi$
              and
              $\Gamma, x:t \vDash E[\top]$.
              %
              From this last one, we can use V-$\forall$i to
              obtain $\Gamma \vDash \forall (x:t). E[\top]$, and conclude
              (noting that $\gamma(P) = x:t, \gamma(P')$).
    \end{itemize}
\end{proof}


\newpage

\section{Modelling native plugins}
\label{app:native}

This appendix contains the formal definitions of our multi-language
interoperability model. We formalize the semantics of source and target language
as well as the translation between the two. We convey the essential ideas 
in the text below, and present full details in 
Figures \ref{fig:sourcelanglong}--\ref{fig:translationlong} that follow the discussion below.

\subsection*{Part 1: Modeling Native Plugins with Simple Types}

%We now present our design of native plugins using a small
%multi-language calculus, a model of interactions between the
%normalizers and native plugins, with appropriate conversions
%between the term representations. We intend to only convey the essential ideas
%of our implementation. A full formalization is out of scope, although more details are in the supplementary material.

\iffalse
\begin{figure}[h]
  \centering
\[\begin{array}{@{}l@{~~}l@{~}r@{}l@{}}
\mbox{Types} & \stau & \bnfdef & \stunit% \bnfalt \svar{A}
      \bnfalt \sarrow{\stau_{\source{1}}}{\stau_{\source{2}}}
      \bnfalt \sprod{\stau_{\source{1}}}{\stau_{\source{2}}} \bnfalt \ssum{\stau_{\source{1}}}{\stau_{\source{2}}} \\
& & & \\

\mbox{Expressions} & \sexp & \bnfdef & \sunit \bnfalt \svar{x} \bnfalt \sabs{\svar{x}}{\stau}{\sexp} \bnfalt \sapp{\sexp_{\source{1}}}{\sexp_{\source{2}}} \\
%& & \bnfalt & \stapp{\sexp}{\stau}\\
& & \bnfalt & \spair{\sexp_{\source{1}}}{\sexp_{\source{2}}} \bnfalt \sfst{\sexp} \bnfalt \ssnd{\sexp} \\
& & \bnfalt & \sinl{\sexp} \bnfalt \sinr{\sexp} \\
& & \bnfalt & \scasex{\sexp}{\stau_{\source{1}}}{\sexp_{\source{1}}}{\stau_{\source{2}}}{\sexp_{\source{2}}} \\
& & \bnfalt & \salien{\texp}{\stau}{\sdelta} \\
% \bnfalt \sopaque{\texp}
\end{array}\]
  \caption{Syntax of the source language}
  \label{fig:syntax-source}
\end{figure}

\begin{figure}[h]
  \centering
\[\begin{array}{@{}l@{~~}l@{~}r@{}l@{}}
\mbox{Expressions} & \texp & \bnfdef & \tunit \bnfalt \tvar{x} \bnfalt \tabs{\tvar{x}}{\texp} \bnfalt \tapp{\texp_{\target{1}}}{\texp_{\target{2}}} \\
& & \bnfalt & \tpair{\texp_{\target{1}}}{\texp_{\target{2}}} \bnfalt \tfst{\texp} \bnfalt \tsnd{\texp} \\
& & \bnfalt & \tinl{\texp} \bnfalt \tinr{\texp} \\
& & \bnfalt & \tcasex{\texp}{\texp_{\target{1}}}{\texp_{\target{2}}} \\
& & \bnfalt & \talien{\sexp}{\stau}{\sdelta} \bnfalt
%\topaque{\sexp}\\

\end{array}\]
  \caption{Syntax of the target language}
  \label{fig:syntax-target}
\end{figure}
\fi

Reflecting our requirement that plugins are only supported at
ML-typeable interfaces, we start with a
source language that is an intrinsically typed (Church
style), standard, simply typed lambda calculus with pairs and
sums. Later in the section, we
will add ML-style rank 1 polymorphism to it.
%
Conversely, reflecting
the type-less native representation of compiled OCaml code, our target
language is the untyped lamdba calculus.
For clarity, we markup the the syntax using colors, using
$\source{blue}$ for the source language and $\target{red}$ for the target.
Aside from the standard forms, we
have two new expression forms to account for the ``alien''
expressions of one language appearing in the other:
$\talien{\sexp}{\stau}{}$, which \emph{embeds} a
$\stau$-typed source language expression into the target language, and 
$\salien{\texp}{\stau}{}$,
which \emph{unembeds} a target language expression to the source language at
type $\stau$.

%% $\salien{\texp}{\stau}{}$,
%% which \emph{unembeds}
%% a target computation in the source language at type $\stau$, and
%% $\talien{\sexp}{\stau}{}$, which
%% \emph{embeds} a $\stau$-typed source computation in the target language.

%% To model native plugins, consider a simply-typed lambda calculus
%% source term $\source{E(\sexp : \stau)}$, a context $\source{E}$
%% surrounding a term $\sexp : \stau$.
%% %
%% % JP: I think these two were salien as opposed to talien?
%% We translate $\sexp$ to an untyped lambda term $\target{\texp}$ (by
%% erasing type information) and unembed it in the original source
%% context as $\source{E (\salien{\target{\texp}}{\stau}{})}$.
%% %
%% We study the interactions of such hybrid terms, aiming to ensure that
%% $\source{E (\salien{\target{\texp}}{\stau}{})}$ is observationally
%% equivalent to $\source{E(\sexp : \stau)}$.

We explain the reductions through a small example. Consider a source language term
$\sapp{\source{(}\sabs{\svar{x}}{\mathbb{\source{Z}}}{\svar{x}}\source{)}}{\source{0}}$, where we want to compile the
identity function to native code and then perform the application. The
first step is to translate the function to the target, by systematically
erasing the types, and unembed it in the source. In our example, the
source language term then becomes
$\sapp{\salien{\tabs{\tvar{x}}{\tvar{x}}}{\mathbb{\source{Z}} \source{->} \mathbb{\source{Z}}}{}}{\source{0}}$.


%% The main characteristic of the rules is that whenever a target term blocks
%% reduction in the source language, we $\sforce{}$ the term into a source value of
%% the appropriate type, using the unembedding function $\scoerce\stau\relax\texp$.
%% When $\sforce{\svar}$ is applied at function type, this involves passing a
%% source term as an argument to a target language function,
%% after \emph{unembedding} said function as a source term, i.e., although not
%% present initially, unembeddings of target terms in the source arise as during
%% reductions.
%% \guido{First use of $\scoerce\stau\relax\texp$, confused. Last sentence
%% seems explained backwards to me, hard to follow.}

%% \jonathan{I'm presenting both reduction semantics in this paragraph, then
%% jumping to the description of the two coercion functions.}

%% \nik{I suggest presenting embeddings first, since
%% this is the direction with which things ``start'' in the plugin
%% scenario.}

%% \zp{What about this structure:
%% \begin{itemize}
%% \item Embeddings : syntax + semantics (reduction and coersions)
%% \item Unembeddings : ditto
%% \item Other (non-standard) rules of the semantics + force
%% \end{itemize}
%% }

% \nik{remove $\delta$ at first and the protection function at first.}
%
%% \begin{mathpar}
%%  \inferrule*[right=Alien]
%%             {\sexp \sevals \sval}
%%             {\talien{\sexp}{\stau}{\relax} \teval \tcoerce{\stau}{\relax}{\sval}}
%% \end{mathpar}

\paragraph{\bf Source semantics and unembedding (\autoref{fig:sourceseman})}
At a beta-reduction step (e.g. rule {\sc{S-App}} below), the source
semantics uses a meta-function
$\sforce{\source{v}}$ to examine the head of the application $\svar{v}$.
In addition to the usual values, the term
$\salien{\tval}{\stau}{\relax}$ is a source value \emph{iff}
$\tval$ is a target value. If the head is such a
value, $\sforce{}$ invokes the unembedding coercion
$\scoerce\stau\relax\tval$ to coerce it to a suitable source
value. For structural types, this amounts to lazily coercing the head
constructor of the term (see below). For function types, a source-level closure is
allocated, which first \emph{embeds} its argument in the target,
reduces a target application (using rule {\sc{S-Alien}}), and
then \emph{unembeds} the result back in the source.
%
%
%% $\sforce{}$ examines its argument and invokes the the embedding meta-function
%% $\scoerce\stau\relax\tval$ whenever it hits an alien value. The embedding
%% meta-function $\scoerce\stau\relax\tval$ takes a target \emph{value} (the target
%% is CBV) and coerces it suitably as a source \emph{value}.
%% \zp{The result of the coercion will always be a value in the simply typed world.
%% It can be an expression when we add poly types.}
%
%
%
%% To define the semantics of the two languages we need a way to interpret alien terms
%% in both contexts. 
%%  The source reduces an alien term $\salien{\texp}{\stau}{\relax}$
%% by invoking the semantics of the target language to obtain a target value (rule {\sc Alien}).
%% Hence, the term $\salien{\tval}{\stau}{\relax}$ is a value in the source \emph{iff}
%% $\tval$ is a value in the target language. Whenever a reduction step is blocked
%% because of an alien value, the semantics will coerce this value to a source value by a
%% call to the meta-function $\sforce{}$.
%
\begin{figure}
\begin{mathpar}
  \inferrule*[right=S-Alien]
             {\texp \tevals \tval}
             {\salien{\texp}{\stau}{} \seval \salien{\tval}{\stau}{}}
  %
%  \and
%  \inferrule*[right=Proj]
%             {\sforce{\sval}=\spair{\sexp_{\source{1}}}{\sexp_{\source{2}}}}
%             {\sval.\source{i} \seval \sexp_{\source{i}}}
 %
  \and
  %
  \inferrule*[right=S-App]
             {\sforce{\sval}=\sabs{\svar{x}}{\stau}{\sexp}}
             {\sapp{\sval}{\sexpii} \seval \sexp \ssubst{\svar{x}}{\sexpii}}
 %
 \and
  \inferrule*{}{
    \sforce{\sval} =
    \begin{cases}
      \scoerce{\stau}{\relax}{\tval}
      & \text{if }\sval = \salien{\tval}{\stau}{\relax} \cr
        \sval
      & \text{otherwise}
    \end{cases}
  }
\end{mathpar}
%
\vspace{0.05cm}
\[
\begin{array}{r l l}
  \scoerce{\stunit}{}{\tunit} & = & \sunit

\\[0.3em]

  \scoerce{\sprod{\staui}{\stauii}}{\relax}{\tpair{\tvali}{\tvalii}} & = &
  \spair{\salien{\tvali}{\staui}{\relax}}{\salien{\tvalii}{\stauii}{\relax}}

\\[0.3em]

  \scoerce{\sarrow{\staui}{\stauii}}{\relax}{\texp} & = &
  \sabs{\svar{x}}{\staui}{\salien{\tapp{\texp}{\talien{\svar{x}}\staui\relax}}{\stauii}{\relax}}
\end{array}
\]
    \caption{Source semantics}
    \label{fig:sourceseman}
\end{figure}

In our example, reduction proceeds by rule {\sc{S-App}} by first
applying the unembedding coercion
$\scoerce{\source{\mathbb{Z} -> \mathbb{Z}}}{}{}$ to $\tabs{\tvar{x}}{\tvar{x}}$ to get
$\sabs{\svar{x}}{\mathbb{\source{Z}}}{\salien{\tapp{\target{(}\tabs{\tvar{x}}{\tvar{x}}\target{)}}{\talien{\svar{x}}{\mathbb{\source{Z}}}{}}}{\mathbb{\source{Z}}}{}}$,
and then applying the substitution to get
$\salien{\tapp{\target{(}\tabs{\tvar{x}}{\tvar{x}}\target{)}}{\talien{\source{0}}{\mathbb{\source{Z}}}{}}}{\mathbb{\source{Z}}}{}$.

%
%\nik{describe embeddings for pairs and then for functions, showing how the rule requires unembedding source terms and passing it to the target language.}
%
%% \jonathan{Probably need to change the rule above to account for the descrepancy
%% between the lambda and the f}
%% As such, when a function $\svar{f}$ is compiled as a plugin, one of
%% \fstar's normalizers reduces
%% $\sapp{\salien{\tvar{f}}{\sarrow\staui\stauii}\relax}{\sexp}$. The source
%% semantics \emph{force} the head of the application by applying the arrow
%% embedding above, followed by rule \textsc{App}, which gives 
%% $
%% \salien{\tapp{\tvar{f}}{\talien\sexp\staui\relax}}{\stauii}{\relax}
%% $.
%% \jonathan{Not sure about this bit.}
%% \zp{From the formalization prospective seems accurate. That about implementation?
%% Does this reflect what happens?}
%% %
%% By virtue of being call by value, the \textsc{Alien} rule for the target
%% language then gives
%% $
%% \salien{\tapp{\tvar{f}}{\tcoerce\staui\relax{\svar{v}}}}{\stauii}{\relax}
%% $.
%% In practice, reducing source terms from within the target language via rule
%% \textsc{Alien} is implemented via a callback to \fstar's normalizer.

\paragraph{\bf Target semantics and embeddings (\autoref{fig:targetseman})}
%% Although not present initially at the top-level, unembeddings of source terms in
%%  the target arise during reductions of embdeddings. Hence, we need to coerce
%%  them in the target.
The target semantics is standard
 CBV.
 %% Embeddings of source terms are not values in the target language. As such,
 %% embedded source terms are fully evaluated when passed to the target language.
 The only new rule is \textsc{T-Alien}, which first reduces the foreign
 source term using  the source semantics, and then proceeds
 to \emph{embed} the resulting value into the target
 language, using the embedding coercion $\tcoerce\stau\relax\sval$,
 as shown below. In our implementation, this rule is realized via a callback
 to the \fstar's normalizer.

 \begin{figure}
\begin{mathpar}
\inferrule*[right=T-Alien]
            {\sexp \sevals \sval}
            {\talien{\sexp}{\stau}{} \teval \tcoerce{\stau}{}{\sval}}
\and
\inferrule*[right=T-App]
            {\quad}
            {\tapp{\target{(}\tabs{\tvar{x}}{\texp}\target{)}}{\tval} \teval \texp\tsubst{\tvar{x}}{\tval}}
\end{mathpar}
%% \jonathan{I'm not sure about having expressions in these rules below, because I
%% thought that the target Alien rule forced reduction of the source expression as
%% a value? Is this meant to imply that we only need to reduce to weak head normal
%% form?}
%% \zp{Yes, we only need whnf to apply the coercions. I'm also using it to indicate that the semantics of the source
%% (i.e. the nornalizer) are CBN and therefore values are in weak head normal form. The rules below look good.
%% Only changed the last one to stress that the inner red term will be a value. }
%Coercions are defined in a similar manner.
\[
\begin{array}{r l l}
  \tcoerce{\sarrow{\staui}{\stauii}}{\relax}\sexp  & = & 
  \tabs{\tvar{x}}{\talien{\sapp\sexp{\salien{\tvar{x}}{\staui}{\relax}}}{\stauii}{\relax}}
  \\[0.3em]
  \tcoerce{\sprod{\staui}{\stauii}}{\relax}{\spair{\sexpi}{\sexpii}}  & = & 
  \tpair{\talien{\sexpi}{\staui}{\relax}}{\talien{\sexpii}{\stauii}{\relax}}
  %
  \\[0.3em] 
  \tcoerce{\staup}{\relax}{\salien{\tval}{\stau}{\relax}}  & = &  \tval
  \text{ if } \stau \equiv \staup
\end{array}
\]
     \caption{Target semantics}
     \label{fig:targetseman}
 \end{figure}

The last case in the definition of the coercion cancels
 superfluous embeddings, in case the term inside is an
 unembedding. Continuing with our example,
$\salien{\tapp{\target{(}\tabs{\tvar{x}}{\tvar{x}}\target{)}}{\talien{\source{0}}{\mathbb{\source{Z}}}{}}}{\mathbb{\source{Z}}}{}$
 reduces using {\sc{S-Alien}}, which reduces the application in its
 premise by first reducing
 $\talien{\source{0}}{\mathbb{\source{Z}}}{}$
 to
$\target{0}$ (using rule {\sc{T-Alien}}, with
 $\tcoerce{\mathbb{Z}}{\relax}{\source{0}} = \target{0}$), and
then using rule {\sc{T-App}} to get
 $\target{0}$. The rule {\sc{S-Alien}} then returns
 $\salien{\target{0}}{\mathbb{\source{Z}}}{}$
 to the \fstar normalizer, which applies the unembedding coercion to
 get $\source{0}$.
 
\iffull

\paragraph{\bf Translation from source to target}
Fig.~\ref{fig:translation} shows a few
selected cases in the translation of source to target terms.
The key rule is \textsc{Abs}, which translates a source
abstraction to a target abstraction. Rather than maintaining an explicit mapping of
source variables to target variables, the translation \emph{embeds} the target
variable $\tvar x$
in the source body of the abstraction. Later on, the \textsc{Box} rule finds the
embedding, and removes it, leaving only the intended variable in the translated
target program. As such, any closed source term is translated to a closed target
term, without any embeddings or unembeddings in it.

Rule \textsc{Var} is therefore only triggered for open variables that have not
been unembedded by \textsc{Abs} already. The formalization accounts
for this, and remains valid in the presence of open
terms. From an implementation perspective, however, we cannot in general compile
open terms to OCaml. Therefore, we restrict ourselves to compiling top-level
definitions, in which the only open variables are other top-level definitions.
Such top-level definitions have a name, which we know how to compile to OCaml.
Because of this restriction, plugin extraction works at the level of top-level
definitions, which are the only AST nodes that can be marked with the
\ls$plugin$ attribute.

\begin{figure}[h]
  \centering
\begin{mathpar}
  \inferrule*[right=Var]
             {\\}
             {\svar{x} \trans \talien{\svar{x}}{\senv\spar{\svar{x}}}{\cdot}}
  %
  \and
  %
  \inferrule*[right=Box]
             {\\}
             {\salien{\texp}{\stau}{\sdelta} \trans \texp}
  %
  \and
  %
  \inferrule*[right=Pair]
             {\source{e_i} \trans \target{e_i}}
             { \spair{\source{e_1}}{\source{e_2}} \trans \tpair{\target{e_1}}{\target{e_2}}}
  %
  \and
  %
  \inferrule*[right=Abs]
             {\sexp\ssubst{\svar{x}}{\salien{\tvar{x}}{\source{\tau}}{\cdot}} \trans \texp}
             {\sabs{\svar{x}}{\stau}{\sexp} \trans \tabs{\tvar{x}}{\texp}}
  %
  \and
  %
  \inferrule*[right=App]
             {\source{e_i} \trans \target{e_i}}
             {\sapp{\source{e_1}}{\source{e_2}} \trans \tapp{\target{e_1}}{\target{e_2}}}
  %
           \end{mathpar}
           \caption{Translation from source to target}
           \label{fig:translation}
\end{figure}
\fi

%% which is then the whole expression reduces to , which finally
%% reduces to  using rule .\guido{\textsc{S-Alien} only seems to
%% return unembedded terms? missing a coercion maybe?}
%% %% \jonathan{Re the rule for arrow types: changed it to no longer do a built-in
%% substitution. (Simpler?)}

%% The base and structural cases (e.g. pairs) are straightforward. The arrow case
%% is symmetrical to $\scoerce{\sarrow\staui\stauii}\relax\relax$. At run-time,
%% this amount to allocating a closure in the target language (OCaml) that calls
%% back into the normalizer, then embeds the result back in the target.
%% \jonathan{I feel like we should say something else in this paragraph, because
%% looking at the definition of force, it seems like we're always applying the
%% cancellation rule (may confuse the reader).}

% ----- Aseem: Begin: Text before starting tilde experiment -----

%% We start with a source language that is an intrinsically typed (Church
%% style), standard, simply-typed lambda calculus with pairs and sums and
%% a target language that is the untyped lambda calculus.
%% %
%% For clarity, we markup the the syntax using colours, using
%% \source{blue} for the source language and \target{red} for the target.
%% %
%% Two constructs are new: $\salien{\texp}{\stau}{}$, which \emph{unembeds}
%% a target computation in the source language at type $\stau$, and
%% $\talien{\sexp}{\stau}{}$, which
%% \emph{embeds} a $\stau$-typed source computation in the target language.

%% To model native plugins, consider a simply-typed lambda calculus
%% source term $\source{E(\sexp : \stau)}$, a context $\source{E}$
%% surrounding a term $\sexp : \stau$.
%% %
%% % JP: I think these two were salien as opposed to talien?
%% We translate $\sexp$ to an untyped lambda term $\target{\texp}$ (by
%% erasing type information) and unembed it in the original source
%% context as $\source{E (\salien{\target{\texp}}{\stau}{})}$.
%% %
%% We study the interactions of such hybrid terms, aiming to ensure that
%% $\source{E (\salien{\target{\texp}}{\stau}{})}$ is observationally
%% equivalent to $\source{E(\sexp : \stau)}$.

%% \iffull

%% Figure \ref{fig:translation} shows a few
%% selected cases in the translation of source to target terms.
%% The key rule is \textsc{Abs}, which translates a source
%% abstraction to a target abstraction. Rather than maintain an explicit mapping of
%% source variables to target variables, the translation \emph{embeds} the target
%% variable $\tvar x$
%% in the source body of the abstraction. Later on, the \textsc{Box} rule finds the
%% embedding, and removes it, leaving only the intended variable in the translated
%% target program. As such, any closed source term is translated to a closed target
%% term, without any embeddings or unembeddings in it.

%% Rule \textsc{Var} is therefore only triggered for open variables that have not
%% been unembedded by \textsc{Abs} already. The formalization accounts
%% for this, and remains valid in the presence of open
%% terms. From an implementation perspective, however, we cannot in general compile
%% open terms to OCaml. Therefore, we restrict ourselves to compiling top-level
%% definitions, in which the only open variables are other top-level definitions.
%% Such top-level definitions have a name, which we know how to compile to OCaml.
%% Because of this restriction, plugin extraction works at the level of top-level
%% definitions, which are the only AST nodes that can be marked with the
%% \ls$plugin$ attribute.

%% \begin{figure}[h]
%%   \centering
%% \begin{mathpar}
%%   \inferrule*[right=Var]
%%              {\\}
%%              {\svar{x} \trans \talien{\svar{x}}{\senv\spar{\svar{x}}}{\cdot}}
%%   %
%%   \and
%%   %
%%   \inferrule*[right=Box]
%%              {\\}
%%              {\salien{\texp}{\stau}{\sdelta} \trans \texp}
%%   %
%%   \and
%%   %
%%   \inferrule*[right=Pair]
%%              {\source{e_i} \trans \target{e_i}}
%%              { \spair{\source{e_1}}{\source{e_2}} \trans \tpair{\target{e_1}}{\target{e_2}}}
%%   %
%%   \and
%%   %
%%   \inferrule*[right=Abs]
%%              {\sexp\ssubst{\svar{x}}{\salien{\tvar{x}}{\source{\tau}}{\cdot}} \trans \texp}
%%              {\sabs{\svar{x}}{\stau}{\sexp} \trans \tabs{\tvar{x}}{\texp}}
%%   %
%%   \and
%%   %
%%   \inferrule*[right=App]
%%              {\source{e_i} \trans \target{e_i}}
%%              {\sapp{\source{e_1}}{\source{e_2}} \trans \tapp{\target{e_1}}{\target{e_2}}}
%%   %
%%            \end{mathpar}
%%            \caption{Translation from source to target}
%%            \label{fig:translation}
%% \end{figure}

%% \fi

%% %% The main characteristic of the rules is that whenever a target term blocks
%% %% reduction in the source language, we $\sforce{}$ the term into a source value of
%% %% the appropriate type, using the unembedding function $\scoerce\stau\relax\texp$.
%% %% When $\sforce{\svar}$ is applied at function type, this involves passing a
%% %% source term as an argument to a target language function,
%% %% after \emph{unembedding} said function as a source term, i.e., although not
%% %% present initially, unembeddings of target terms in the source arise as during
%% %% reductions.
%% %% \guido{First use of $\scoerce\stau\relax\texp$, confused. Last sentence
%% %% seems explained backwards to me, hard to follow.}

%% %% \jonathan{I'm presenting both reduction semantics in this paragraph, then
%% %% jumping to the description of the two coercion functions.}

%% %% \nik{I suggest presenting embeddings first, since
%% %% this is the direction with which things ``start'' in the plugin
%% %% scenario.}

%% %% \zp{What about this structure:
%% %% \begin{itemize}
%% %% \item Embeddings : syntax + semantics (reduction and coersions)
%% %% \item Unembeddings : ditto
%% %% \item Other (non-standard) rules of the semantics + force
%% %% \end{itemize}
%% %% }

%% % \nik{remove $\delta$ at first and the protection function at first.}
%% %
%% %% \begin{mathpar}
%% %%  \inferrule*[right=Alien]
%% %%             {\sexp \sevals \sval}
%% %%             {\talien{\sexp}{\stau}{\relax} \teval \tcoerce{\stau}{\relax}{\sval}}
%% %% \end{mathpar}

%% \paragraph{Embeddings}
%% To define the semantics of the two languages we need a way to interpret alien terms
%% in both contexts. 
%%  The source reduces an alien term $\salien{\texp}{\stau}{\relax}$
%% by invoking the semantics of the target language to obtain a target value (rule {\sc Alien}).
%% Hence, the term $\salien{\tval}{\stau}{\relax}$ is a value in the source \emph{iff}
%% $\tval$ is a value in the target language. Whenever a reduction step is blocked
%% because of an alien value, the semantics will coerce this value to a source value by a
%% call to the meta-function $\sforce{}$.
%% %
%% \begin{mathpar}
%%   \inferrule*[right=Alien]
%%              {\texp \tevals \tval}
%%              {\salien{\texp}{\stau}{\sdelta} \seval \scoerce{\stau}{\sdelta}{\tval}}
%%   %
%% %  \and
%% %  \inferrule*[right=Proj]
%% %             {\sforce{\sval}=\spair{\sexp_{\source{1}}}{\sexp_{\source{2}}}}
%% %             {\sval.\source{i} \seval \sexp_{\source{i}}}
%%  %
%%   \and
%%   %
%%   \inferrule*[right=App]
%%              {\sforce{\sval}=\sabs{\svar{x}}{\stau}{\sexp}}
%%              {\sapp{\sval}{\sexpii} \seval \sexp \ssubst{\svar{x}}{\sexpii}}
%%  %
%%  \and
%%   \inferrule*{}{
%%     \sforce{\sval} =
%%     \begin{cases}
%%       \scoerce{\stau}{\relax}{\tval}
%%       & \text{if }\sval = \salien{\tval}{\stau}{\relax} \cr
%%         \sval
%%       & \text{otherwise}
%%     \end{cases}
%%   }
%% \end{mathpar}
%% $\sforce{}$ examines its argument and invokes the the embedding meta-function
%% $\scoerce\stau\relax\tval$ whenever it hits an alien value. The embedding
%% meta-function $\scoerce\stau\relax\tval$ takes a target \emph{value} (the target
%% is CBV) and coerces it suitably as a source \emph{value}.
%% \zp{The result of the coercion will always be a value in the simply typed world.
%% It can be an expression when we add poly types.}
%% For structural types, this amounts to lazily coercing the head constructor of
%% the term (below). For function types, a source-level closure is allocated, which
%% first \emph{embeds} its argument in the target, reduces a target application,
%% then \emph{unembeds} the result back in the source.
%% %
%% \[
%% \begin{array}{c c}
%%   \scoerce{\sprod{\staui}{\stauii}}{\relax}{\tpair{\tvali}{\tvalii}} =
%%   \spair{\salien{\tvali}{\staui}{\relax}}{\salien{\tvalii}{\stauii}{\relax}}

%%   \qquad

%% &

%%   \qquad

%%   \scoerce{\sarrow{\staui}{\stauii}}{\relax}{\tabs{\tvar{x}}{\texp}} =
%%   \sabs{\svar{x}}{\staui}{\salien{\tapp{(\tabs{\tvar{x}}{\texp})}{\talien{\svar{x}}\staui\relax}}{\stauii}{\relax}}
%% \end{array}
%% \]
%% %
%% %\nik{describe embeddings for pairs and then for functions, showing how the rule requires unembedding source terms and passing it to the target language.}
%% %
%% \jonathan{Probably need to change the rule above to account for the descrepancy
%% between the lambda and the f}
%% As such, when a function $\svar{f}$ is compiled as a plugin, one of
%% \fstar's normalizers reduces
%% $\sapp{\salien{\tvar{f}}{\sarrow\staui\stauii}\relax}{\sexp}$. The source
%% semantics \emph{force} the head of the application by applying the arrow
%% embedding above, followed by rule \textsc{App}, which gives 
%% $
%% \salien{\tapp{\tvar{f}}{\talien\sexp\staui\relax}}{\stauii}{\relax}
%% $.
%% \jonathan{Not sure about this bit.}
%% \zp{From the formalization prospective seems accurate. That about implementation?
%% Does this reflect what happens?}
%% %
%% By virtue of being call by value, the \textsc{Alien} rule for the target
%% language then gives
%% $
%% \salien{\tapp{\tvar{f}}{\tcoerce\staui\relax{\svar{v}}}}{\stauii}{\relax}
%% $.

%% \paragraph{Unembeddings}
%% Although not present initially at the top-level, unembeddings of source terms in
%%  the target arise during reductions of embdeddings. Hence, we need to coerce
%%  them in the target. The reduction semantics for the target language is standard
%%  CBV. Embeddings of source terms are not values in the target language. As such,
%%  embedded source terms are fully evaluated when passed to the target language.
%%  The only new rule is \textsc{Alien}, which first reduces the foreign term using
%%  the source language semantics, then proceeds to \emph{embed} it into the target
%%  language, using the target \emph{unembedding} $\tcoerce\stau\relax\sval$.
%% \begin{mathpar}
%% \inferrule*[right=Alien]
%%             {\sexp \sevals \sval}
%%             {\talien{\sexp}{\stau}{\sdelta} \teval \tcoerce{\stau}{\sdelta}{\sval}}
%% \end{mathpar}
%% %% \jonathan{I'm not sure about having expressions in these rules below, because I
%% %% thought that the target Alien rule forced reduction of the source expression as
%% %% a value? Is this meant to imply that we only need to reduce to weak head normal
%% %% form?}
%% %% \zp{Yes, we only need whnf to apply the coercions. I'm also using it to indicate that the semantics of the source
%% %% (i.e. the nornalizer) are CBN and therefore values are in weak head normal form. The rules below look good.
%% %% Only changed the last one to stress that the inner red term will be a value. }
%% Coercions are defined in a similar manner.
%% \[
%% \begin{array}{c c c}
%%   \tcoerce{\sarrow{\staui}{\stauii}}{\relax}\sexp  = 
%%   \tabs{\tvar{x}}{\talien{\sapp\sexp{\salien{\tvar{x}}{\staui}{\relax}}}{\stauii}{\relax}}
%%   & 
%%   \quad
%%   \tcoerce{\sprod{\staui}{\stauii}}{\relax}{\spair{\sexpi}{\sexpii}}  = 
%%   \tpair{\talien{\sexpi}{\staui}{\relax}}{\talien{\sexpii}{\stauii}{\relax}}
%%   %
%%   \quad
%%   & 
%%   \tcoerce{\staup}{\relax}{\salien{\tval}{\stau}{\relax}}  =  \tval
%%   \text{ if } \stau \equiv \staup
%% \end{array}
%% \]
%% \jonathan{Re the rule for arrow types: changed it to no longer do a built-in
%% substitution. (Simpler?)}

%% The base and structural cases (e.g. pairs) are straightforward. The arrow case
%% is symmetrical to $\scoerce{\sarrow\staui\stauii}\relax\relax$. At run-time,
%% this amount to allocating a closure in the target language (OCaml) that calls
%% back into the normalizer, then embeds the result back in the target. The last
%% case covers the cancellation of embeddings and un-embeddings, and avoids
%% generating superfluous unembeddings. 
%% \jonathan{I feel like we should say something else in this paragraph, because
%% looking at the definition of force, it seems like we're always applying the
%% cancellation rule (may confuse the reader).}

% ----- Aseem: End: Text before starting tilde experiment -----

\subsection*{Part 2: Adding ML-style Polymorphism} %\nik{Budget: 1/2 page}}
\label{sec:parametricity}

%% We extend the source language with ML-style polymorphism and show how
%% we use it to support a limited but useful form of open reduction.
%%  This
%% facilitates the use of polymorphic libraries (e.g., classic library
%% functions on lists) as native plugins.

The key observation in adding ML-style polymorphism is as
follows. Since we only consider target terms
compiled from well-typed source terms, by virtue of
parametricity, a compiled polymorphic function must treat its
argument abstractly. As a result, the normalizers can pass such
arguments \emph{as-is} without applying any embedding coercions---passing
an open term is simply a subcase without further
difficulty. Thus, we can leverage polymorphism to support such limited
but useful open reductions.

In the formal model, we add an \emph{opaque} construct to the target
language, written as
$\topaque\sexp$, which denotes an embedding of $\sexp$ at some
abstract type $\source{A}$ (as mentioned above, coercion to opaque is
a no-op in the implementation\guido{not mentioned}).
%
The coercion functions are extended
with rules to introduce and eliminate opaque terms (the $\sdelta$
superscript carries type substitutions as a technical device), e.g.:
\[
\begin{array}{c c c}
  \tcoerce{\svar{A}}{\sdelta}{\sexp}  = \topaque{\sexp}
  \quad & \quad
  \scoerce{\svar{A}}{\sdelta}{\topaque{\sexp}} = \sexp
  \quad
  &
  \quad
  \tcoerce{\sprod{\staui}{\stauii}}{\sdelta}{\spair{\sexpi}{\sexpii}} =
  \tpair{\prot{\source{e_1}}{\source{\tau_1}}{\sdelta}}{\prot{\source{e_2}}{\source{\tau_2}}{\sdelta}}
\end{array}
\]
where $\prot{\sexp}{\stau}{\sdelta}$ is $\topaque{\sexp}$ when
$\stau = \source{A}$, or a usual type-directed embedding otherwise. The source
semantics is 
extended with a rule for type applications, while target
semantics, being untyped, remains as-is.

Consider an application of the polymorphic identity function to a
variable $\svar{y}$:
$\sapp{\stapp{(\stabs{\svar{A}}{\sabs{\svar{x}}{\svar{A}}{\svar{x}}})}{\source{\mathbb{Z}}}}{\svar{y}}$.
As before, we would like to compile the function to native code and
then reduce the applications---except this time it is an open term.
%
Translation of the function and its embedding in the source yields
$\sapp{\stapp{\salien{\tabs{\tvar{x}}{\tvar{x}}}{\sall{\svar{A}}{\sarrow{\svar{A}}{\svar{A}}}}{\sesubstemp}}{\source{\mathbb{Z}}}}{\svar{y}}$.
%
We now apply the unembedding coercion
$\scoerce{\sall{\svar{A}}{\sarrow{\svar{A}}{\svar{A}}}}{\sesubstemp}{\tabs{\tvar{x}}{\tvar{x}}}$
which returns \linebreak a source type abstraction
$\stabs{\svar{A}}{\scoerce{\sarrow{\svar{A}}{\svar{A}}}{\sesubstemp}{\tabs{\tvar{x}}{\tvar{x}}}}$.
%
Reducing the arrow coercion next, we get
$\stabs{\svar{A}}{\sabs{\svar{x}}{\svar{A}}{\salien{\tapp{(\tabs{\tvar{x}}{\tvar{x}})}{\prot{\svar{x}}{\svar{A}}{\sesubstemp}}}{\svar{A}}{\sesubstemp}}}$,
%
where $\prot{\svar{x}}{\svar{A}}{\sesubstemp} = \topaque{\svar{x}}$, as mentioned above.
This is the key idea---when the original application reaches the beta-reduction with
argument $\svar{y}$ (after the type application is reduced), $\svar{y}$ is simply substituted
in this opaque construct, which is then returned back to the \fstar normalizer as
$\salien{\topaque{\svar{y}}}{\svar{A}}{\sesubstsingl{\svar{A}}{\source{\mathbb{Z}}}}$.
%
The normalizer can then apply the identity unembedding and get $\svar{y}$.

This ability to reduce open terms relying on polymorphism allows us to
evaluate expressions like
$\sapp{\sapp{\stapp{\stapp{\salien{\target{\text{\ls{List.map}}}}{\sall{\svar{A}}{\sall{\svar{B}}{\sarrow{\sarrow{(\sarrow{\svar{A}}{\svar{B}})}{\source{\text{\ls{list}}} \svar{A}}}{\source{\text{\ls{list}}} \svar{B}}}}}{\sesubstemp}}{\source{\mathbb{Z}}}}{\source{\mathbb{Z}}}}{\source{\text{\ls{f}}}}}{\source{\text{\ls{[0; ...; 10}$^6$\ls{]}}}}$,
with \ls$f$ a free variable,
using native datastructure implementations (\ls{List} in this case), 
which is much faster than using the normalizers.
% (see \autoref{sec:registers}).

\subsection*{Part 3: Full Details on our Multi-language Interoperability Model}

%% \aseem{TODO: add an example of applying the polymorphic identity
%% function to a variable. Add the property we hope for.}

%% prevents further reduction or coercion (it is a value in
%% the target language). An opaque construct preserves a source term at an abstract type.
%% %
%% We define the abstraction-protecting embedding $\prot\sexp\stau\sdelta$ to be
%% either a regular $\talien\sexp\stau\sdelta$ as before, or an opaque
%% $\topaque\sexp$ if $\sdelta(\stau) = \svar{A}$. Rules that previously generated
%% $\talien\sexp\stau\sdelta$ now generate $\prot\sexp\stau\sdelta$, e.g.
%% $
%%   \tcoerce{\sarrow{\staui}{\stauii}}{\sdelta}\sexp  = 
%%   \tabs{\tvar{x}}{\prot{\sapp\sexp{\salien{\tvar{x}}{\staui}{\sdelta}}}{\stauii}{\sdelta}}
%% $


%% In the formal model, we add a new \em

%% In a polymorphic setting, abstract type variables appear in the language.
%% To ensure soundness, coercions in the target language must stop at such abstract
%% type variables. (Coercing underneath an abstraction would reveal the underlying
%% type to the source language, and return an ill-typed term.) We add a type
%% substitution $\sdelta$ in:
%% $\tcoerce\stau\sdelta\sval$, 
%% $\scoerce\stau\sdelta\tval$, 
%% $\talien\sval\stau\sdelta$, 
%% $\salien\tval\stau\sdelta$. Upon entering a type scheme
%% $\sall{\svar{A}}{\ssig}$, $\sdelta$ is extended and becomes
%% ${\sesubstext{\sdelta}{\svar{A}}{\stau}}$.

%% We extend the target language with an \emph{opaque} construct, denoted
%% $\topaque\sexp$, which prevents further reduction or coercion (it is a value in
%% the target language). An opaque construct preserves a source term at an abstract type.
%% %
%% We define the abstraction-protecting embedding $\prot\sexp\stau\sdelta$ to be
%% either a regular $\talien\sexp\stau\sdelta$ as before, or an opaque
%% $\topaque\sexp$ if $\sdelta(\stau) = \svar{A}$. Rules that previously generated
%% $\talien\sexp\stau\sdelta$ now generate $\prot\sexp\stau\sdelta$, e.g.
%% $
%%   \tcoerce{\sarrow{\staui}{\stauii}}{\sdelta}\sexp  = 
%%   \tabs{\tvar{x}}{\prot{\sapp\sexp{\salien{\tvar{x}}{\staui}{\sdelta}}}{\stauii}{\sdelta}}
%% $

%% Crucially, the coercion functions are extended with rules to deal with opaque
%% terms.
%% \[
%% \begin{array}{c c}
%%   \tcoerce{\svar{A}}{\sdelta}{\sexp}  = \topaque{\sexp}
%%   \quad & \quad
%%   \scoerce{\svar{A}}{\sdelta}{\topaque{\sexp}} = \sexp
%% \end{array}
%% \]
%% In essence, these rules embody the fact that a source function that is
%% parametric over type $\svar{A}$ cannot examine a value at this abstract type,
%% meaning that there actually is no need to unembed or embed such at value at the
%% language boundary: the target program cannot do anything with it, except return
%% it.

%% One consequence of this is that by using polymorphic functions, the programmer
%% can devise a library of natively-compiled plugins that are capable of dealing
%% with open terms. The polymorphic versions are compiled natively, but are used
%% from the source at concrete types that may contain open variables. No embeddings
%% are generated, meaning that the target code is oblivious that the values it is
%% passing around are opaque, open source terms.

%% \jonathan{Is this what we want to say?}
%% Naturally, \fstar goes beyond ML-style polymorphism. For the time being, we
%% decline to compile plugins that, at extraction-time, cannot be re-typed with ML
%% polymorphism.
%% \aseem{We mentioned it before in the limitations of native plugins.}


%\newcommand\path{interop_lazy/}%

\begin{wrapfigure}[htb]{l}{\textwidth}%
    %% \centering
    %% \begin{minipage}{0.54\textwidth}\label{fig:srcexp}
      \centering
      \begin{framed}
        \[\begin{array}{@{}l@{~~}l@{~}r@{}l@{}}
\mbox{Types} & \stau & \bnfdef & \stunit \bnfalt \svar{A}
      \bnfalt \sarrow{\stau_{\source{1}}}{\stau_{\source{2}}}
      \bnfalt \sprod{\stau_{\source{1}}}{\stau_{\source{2}}} \bnfalt \ssum{\stau_{\source{1}}}{\stau_{\source{2}}} \\
& & & \\

\mbox{Type Schemes} & \ssig & \bnfdef & \sall{\svar{A}}{\ssig} \bnfalt \stau \\
& & & \\


\mbox{Environments} & \senv & \bnfdef & \senvemp \bnfalt \senvext{\senv}{\svar{x}}{\stau} \\
& \stenv & \bnfdef & \senvemp \bnfalt \senvtext{\stenv}{\svar{A}} \\
& & & \\
\mbox{Substitutions} & \sesubst & \bnfdef & \sesubstemp \bnfalt \sesubstext{\sesubst}{\svar{A}}{\stau} \\
& & & \\
\mbox{Expressions} & \sexp & \bnfdef & \sunit \bnfalt \svar{x} \bnfalt \sabs{\svar{x}}{\stau}{\sexp} \bnfalt \sapp{\sexp_{\source{1}}}{\sexp_{\source{2}}} \\
& & \bnfalt & \stabs{A}{e} \bnfalt \stapp{\sexp}{\stau} \\
& & \bnfalt & \spair{\sexp_{\source{1}}}{\sexp_{\source{2}}} \bnfalt \sfst{\sexp} \bnfalt \ssnd{\sexp} \\
& & \bnfalt & \sinl{\sexp} \bnfalt \sinr{\sexp} \\
& & \bnfalt & \scasex{\sexp}{\stau_{\source{1}}}{\sexp_{\source{1}}}{\stau_{\source{2}}}{\sexp_{\source{2}}} \\
& & \bnfalt & \salien{\texp}{\stau}{\sdelta} \\
% \bnfalt \sopaque{\texp}
\end{array}\]

      \end{framed}
      \caption{Source Language}
      \label{fig:sourcelanglong}
    %% \end{minipage}\hfill
    %% \begin{minipage}{0.44\textwidth}\label{fig:trgexp}
\end{wrapfigure}


\begin{wrapfigure}[htb]{l}{\textwidth}%
      \centering
      \begin{framed}
        \[\begin{array}{@{}l@{~~}l@{~}r@{}l@{}}
\mbox{Expressions} & \texp & \bnfdef & \tunit \bnfalt \tvar{x} \bnfalt \tabs{\tvar{x}}{\texp} \bnfalt \tapp{\texp_{\target{1}}}{\texp_{\target{2}}} \\
& & \bnfalt & \tpair{\texp_{\target{1}}}{\texp_{\target{2}}} \bnfalt \tfst{\texp} \bnfalt \tsnd{\texp} \\
& & \bnfalt & \tinl{\texp} \bnfalt \tinr{\texp} \\
& & \bnfalt & \tcasex{\texp}{\texp_{\target{1}}}{\texp_{\target{2}}} \\
& & \bnfalt & \talien{\sexp}{\stau}{\sdelta} \bnfalt
\topaque{\sexp}\\
& & & \\
\mbox{Environments} & \tenv & \bnfdef & \tenvemp \bnfalt \tenvext{\tenv}{\tvar{x}} \\

\end{array}\]

      \end{framed}
      \caption{Target Language}
%    \end{minipage}
\end{wrapfigure}

\begin{wrapfigure}[htb]{l}{\textwidth}%
  \label{fig:srctyp}
  \fbox{$\iswf{\stenv}{\stau}$}
  \begin{framed}
    \begin{mathpar}
  \inferrule*[right=TyVar]
  {\\}
  {\iswf{\senvtext{\stenv}{\svar{A}}}{\svar{A}}}
  %
  \and
  %
  \inferrule*[right=Scheme]
  {\iswf{\senvtext{\stenv}{\svar{A}}}{\sexp}}
  {\iswf{\stenv}{\sall{\svar{A}}{\stau}}}
  %
  \and
  %
  \inferrule*[right=Env]
  {\forall ~x,~\iswf{\stenv}{\senv(x)}}
  {\iswf{\stenv}{\senv}}
  %% \inferrule*[right=Opaque]
  %% {\\}
  %% {\hastyp{\senv}{\sopaque{\texp}}{\sttop}}
\end{mathpar}

  \end{framed}
  \caption{Type Well-formedness}
\end{wrapfigure}

\begin{wrapfigure}[htb]{l}{\textwidth}%
  \label{fig:srctyp}
  \fbox{$\iswf{\stenv}{\stau}$}
  \begin{framed}
    \begin{mathpar}
  \inferrule*[right=Empty]
  {\\}
  {\iswf{\cdot}{\cdot}}
  %
  \and
  %
  \inferrule*[right=Cons]
             {
               \iswf{\stenv}{\sdelta} \\
               \iswf{\stenv}{\stau} \\
             }
             {
               \iswf{\senvtext{\stenv}{\svar{A}}}
                    {\sesubstext{\sdelta}{\svar{A}}{\stau}}
             }
\end{mathpar}

  \end{framed}
  \caption{Type Substitution Well-formedness}
\end{wrapfigure}


\begin{wrapfigure}[htb]{l}{\textwidth}%
  \label{fig:srctyp}
  \fbox{$\hastyp{\stenv}{\senv}{\sexp}{\stau}$}
  \begin{framed}
    \begin{mathpar}
\inferrule*[right=Alien]
  {\isok{\stenv}{\senv}{\texp}}
  {\hastyp{\stenv}{\senv}{\salien{\texp}{\stau}{\sesubst}}{\sesubst \stau}}
  %% %
  %% \and
  %% %
  %% \inferrule*[right=Opaque]
  %% {\\}
  %% {\hastyp{\senv}{\sopaque{\texp}}{\sttop}}
\end{mathpar}

  \end{framed}
  \caption{Source Typing}
\end{wrapfigure}

%% \begin{wrapfigure}[htb]{l}{\textwidth}%
%%   \label{fig:srctyp}
%%   \fbox{$\isok{\stenv}{\senv}{\texp}$ (This might not be needed)}
%%   \begin{framed}
%%     \input{trg_typing}
%%   \end{framed}
%%   \caption{Target well-formedness}
%% \end{wrapfigure}

\begin{wrapfigure}[htb]{l}{\textwidth}%
  \label{fig:srceval}
  \fbox{$\sexp \seval \sexpp$}
  \begin{framed}
    \begin{mathpar}
  \inferrule*[right=Ctx]
             {\sexp \seval \sexpp}
             {\scon{\sexp} \seval \scon{\sexpp}}
  %
  \and
  \inferrule*[right=App]
             {\sforce{\sval}=\sabs{\svar{x}}{\stau}{\sexp}}
             {\sapp{\sval}{\sexpii} \seval \sexp \ssubst{\svar{x}}{\sexpii}}

  %
  \and
  %
  \inferrule*[right=TApp]
             {\sforce{\sval}=\stabs{\svar{A}}{\sexpp}}
             { \stapp{\sval}{\stau} \seval
               \sexpp \ssubst{\svar{A}}{\stau}}
  %
  \and
  %
  \inferrule*[right=CaseL]
             {\sforce{\sval} = \sinl{\sexpp}}
             {\scasex{\sval}{\stau_{\source{1}}}{\sexp_{\source{1}}}{\stau_{\source{2}}}{\sexp_{\source{2}}} \seval \sexp_{\source{1}} \ssubst{\svar{x}}{\sexpp}}
  %
  \and
  %
    \inferrule*[right=CaseR]
             {\sforce{\sval} = \sinr{\sexpp}}
             {\scasex{\sval}{\stau_{\source{1}}}{\sexp_{\source{1}}}{\stau_{\source{2}}}{\sexp_{\source{2}}} \seval \sexp_{\source{2}} \ssubst{\svar{x}}{\sexpp}}
  %
  \and
  %
  \inferrule*[right=Proj]
             {\sforce{\sval}=\spair{\sexp_{\source{1}}}{\sexp_{\source{2}}}}
             {\sval.\source{i} \seval \sexp_{\source{i}}}
 %
  \and
  %
  \inferrule*[right=Alien]
             {\texp \tevals \tval}
             {\salien{\texp}{\stau}{\sdelta} \seval \salien{\sval}{\stau}{\sdelta}}
   %
  \and
  %
  %% \inferrule*[right=ProjAlien]
  %%            {\\}
  %%            {\salien{\tpair{\texp_{\target{1}}}{\texp_{\target{2}}}}{\sprod{\stau_{\source{1}}}{\stau_{\source{2}}}}{\sdelta}
  %%              \seval
  %%              \salien{\texp_{\target{1}}}{\stau_{\source{1}}}{\sdelta} }
  %% %
  %% \and
  %% %
  %% \inferrule*[right=CaseLAlien]
  %%            {\\}
  %%            {\scasex{\salien{\tinl{\texp}}{\ssum{\stau_{\source{1}}}{\stau_{\source{2}}}}{\sdelta}}{\stau_{\source{1}}}{\sexp_{\source{1}}}{\stau_{\source{2}}}{\sexp_{\source{2}}}
  %%              \seval
  %%              \sexp_{\source{1}} \ssubst{\svar{x}}{\salien{\texp}{\stau_{\source{1}}}{\sdelta}}}
  %% %
  %% \and
  %% %
  %% \inferrule*[right=AppAlien]
  %%            {\\}
  %%            {\sapp{\salien{\tabs{\tvar{x}}{\texp}}{\sarrow{\stau_{\source{1}}}{\stau_{\source{2}}}}{\sdelta}}
  %%              {\sexp}
  %%              \seval
  %%              %% \sabs{\svar{x}}{\sdelta\stau_{\source{1}}}
  %%                   {\kap{\stau_{\source{1}}}{\sdelta}{\sexp}
  %%                     {\mfun{\texpp}{\salien{\texp\ssubst{\tvar{x}}{\texpp}}{\stau_{\source{2}}}{\sdelta}}}}
  %%            }
  %% %
  %% \and
  %% %
  %% \inferrule*[right=TAppAlien]
  %%            {\\}
  %%            {\stapp{\salien{\texp}{\sall{\svar{A}}{\stau}}{\sdelta}}{\staup}
  %%              \seval
  %%              \salien{\texp}{\stau}{\sesubstext{\sdelta}{\svar{A}}{\staup}}
  %%            }
  %% %
  %% \and
  %% %
  %% \inferrule*[right=AlienOpaque]
  %%            {\stau \not = \sttop}
  %%            {\salien{\topaque{\sexp}}{\stau}{\sdelta}
  %%              \seval
  %%              \sexp
  %%            }
  %% %
  %% \and
  %% %
  %% \inferrule*[right=AlienCancel]
  %%            {\\}
  %%            {\salien{\talien{\sexp}{\staup}{\sdeltap}}{\stau}{\sdelta}
  %%              \seval
  %%              \sexp}
\end{mathpar}

    where
    \[\begin{array}{r l l}
  \sforce{\sval} =
  \begin{cases}
    \scoerce{\stau}{\sdelta}{\tval}
    & \text{if }\sexp = \salien{\tval}{\stau}{\sdelta} \text{ and }
      \sexp \sevals \sval
\\
      \sval
    & \text{otherwise}
  \end{cases}
\end{array}\]

  \end{framed}
  \caption{Evaluation in the Source Language}
\end{wrapfigure}

\begin{wrapfigure}[htb]{l}{\textwidth}%
  \label{fig:trgeval}
  \fbox{$\texp \teval \texpp$}
  \begin{framed}
    \begin{mathpar}
    \inferrule*[right=Ctx]
             {\texp \teval \texpp}
             {\tcon{\texp} \teval \tcon{\texpp}}
  %
  \and
  %
  \inferrule*[right=App]
             {\\}
             {\tapp{(\tabs{\tvar{x}}{\texp})}{\tval} \teval \texp \tsubst{\tvar{x}}{\tval}}

  %
  \and
  %
  \inferrule*[right=CaseL]
             {\\}
             {\tcasex{\tinl{\tval}}{\texp_{\target{1}}}{\texp_{\target{2}}} \teval \texp_{\target{1}} \tsubst{\tvar{x}}{\tval}}
  %
  \and
  %
  \inferrule*[right=CaseL]
             {\\}
             {\tcasex{\tinr{\tval}}{\texp_{\target{1}}}{\texp_{\target{2}}} \teval \texp_{\target{2}} \tsubst{\tvar{x}}{\tval}}
  %
  \and
  %
  \inferrule*[right=Proj]
             {\\}
             {\tpair{\tval_{\target{1}}}{\tval_{\target{2}}}\target{.i} \teval \tval_{\target{i}}}
 %
 \and
 %
 \inferrule*[right=Alien]
            {\sexp \sevals \sval}
            {\talien{\sexp}{\stau}{\sdelta} \teval \tcoerce{\stau}{\sdelta}{\sval}}
 %
 \and
 %%  %
 %%  \inferrule*[right=ProjAlien]
 %%             {\\}
 %%             {\salien{\tpair{\texp_{\target{1}}}{\texp_{\target{2}}}}{\sprod{\stau_{\source{1}}}{\stau_{\source{2}}}}{\sdelta}
 %%               \seval
 %%               \salien{\texp_{\target{1}}}{\stau_{\source{1}}}{\sdelta} }
 %%  %
 %%  \and
 %%  %
 %%  \inferrule*[right=CaseLAlien]
 %%             {\\}
 %%             {\scasex{\salien{\tinl{\texp}}{\ssum{\stau_{\source{1}}}{\stau_{\source{2}}}}{\sdelta}}{\stau_{\source{1}}}{\sexp_{\source{1}}}{\stau_{\source{2}}}{\sexp_{\source{2}}}
 %%               \seval
 %%               \sexp_{\source{1}} \ssubst{\svar{x}}{\salien{\texp}{\stau_{\source{1}}}{\sdelta}}}
 %%  %
 %%  \and
 %%  %
 %%  \inferrule*[right=AppAlien]
 %%             {\\}
 %%             {\sapp{\salien{\tabs{\tvar{x}}{\texp}}{\sarrow{\stau_{\source{1}}}{\stau_{\source{2}}}}{\sdelta}}
 %%               {\sexp}
 %%               \seval
 %%               \sabs{\svar{x}}{\sdelta\stau_{\source{1}}}
 %%                    {\kap{\stau_{\source{1}}}{\sdelta}{\svar{x}}
 %%                      {\mfun{\texpp}{\salien{\texp\ssubst{\tvar{x}}{\texpp}}{\stau_{\source{2}}}{\sdelta}}}}
 %%             }
 %%  %
 %%  \and
 %%  %
 %%  \inferrule*[right=TAppAlien]
 %%             {\\}
 %%             {\stapp{\salien{\texp}{\sall{\svar{A}}{\stau}}{\sdelta}}{\staup}
 %%               \seval
 %%               \salien{\texp}{\stau}{\sesubstext{\sdelta}{\svar{A}}{\staup}}
 %%             }
 %%  %
 %%  \and
 %%  %
 %%    \inferrule*[right=AlienOpaque]
 %%             {\\}
 %%             {\salien{\topaque{\sexp}}{\stau}{\sdelta}
 %%               \seval
 %%               \sexp
 %%             }
 %%               %
 %%  \and
 %%  %
  %% \inferrule*[right=AlienCancel]
  %%            {\\}
  %%            {\talien{\salien{\texp}{\staup}{\sdeltap}}{\stau}{\sdelta}
  %%              \teval
  %%              \texp}
\end{mathpar}

  \end{framed}
  \caption{Evaluation in the Target Language}
\end{wrapfigure}

\medskip

\begin{wrapfigure}[htb]{l}{\textwidth}%
    \centering
    \begin{minipage}[t]{0.45\textwidth}\label{fig:stypsubst}
      \centering
      \begin{framed}
        \[
\begin{array}{r l l}
  \sunit \ssubst{\svar{A}}{\stau} & = & \sunit
  \\[0.3em]
  %
  \sabs{\svar{x}}{\staup}{\sexp} \ssubst{\svar{A}}{\stau} & = &
  \sabs{\svar{x}}{\staup\ssubst{\svar{A}}{\stau}}{\sexp\ssubst{\svar{A}}{\stau}}
  \\[0.3em]
  %
  \vdots 
  \\[0.3em]
  \salien{\texp}{\staup}{\sdelta}\ssubst{\svar{A}}{\stau}  & = &
  \salien{\texp\ssubst{\svar{A}}{\stau}}{\staup}
 {\sesubstext{\sdelta}{\svar{A}}{\stau}}
  \\[0.3em]
  %
  %
  %
  %% \scoerce{\sttop}{\sdelta}{\salien{\tval}{\staup}{\sdeltap}} & = &
  %% \tval
  %% \\[0.3em]
\end{array}
\]

      \end{framed}
      \caption{Type Substitution in Source Terms}
    \end{minipage}\hfill
    \begin{minipage}[t]{0.45\textwidth}\label{fig:ttypsubst}
      \centering
      \begin{framed}
       \[
\begin{array}{r l l}
  \tunit \ssubst{\svar{A}}{\stau} & = & \tunit
  \\[0.3em]
  %
  \tabs{\tvar{x}}{\texp} \ssubst{\svar{A}}{\stau} & = &
  \tabs{\tvar{x}}{\texp\ssubst{\svar{A}}{\stau}}
  \\[0.3em]
  %
  \vdots 
  \\[0.3em]
  \talien{\sexp}{\staup}{\sdelta}\ssubst{\svar{A}}{\stau}  & = &
  \talien{\sexp\ssubst{\svar{A}}{\stau}}{\staup}{\sesubstext{\sdelta}{\svar{A}}{\stau}}
  \\[0.3em]
  \topaque{\sval}  & = & 
  \topaque{\sval \ssubst{\svar{A}}{\stau}}
  \\[0.3em]
\end{array}
\]

      \end{framed}
      \caption{Type Substitution in Target Terms}
    \end{minipage}
\end{wrapfigure}


\medskip

\begin{wrapfigure}[htb]{L}{\textwidth}%
  \label{fig:trans}
  \fbox{$\sexp \trans \texp$} \\
  $\senv$ is a global parameter 
  \begin{framed}
    \begin{mathpar}
  \inferrule*[right=Unit]
             {\\}
             {\sunit \trans \tunit}
  %
  \and
  \inferrule*[right=Var]
             {\\}
             {\svar{x} \trans \talien{\svar{x}}{\senv\spar{\svar{x}}}{\cdot}}

  %
  \and
  %
  \inferrule*[right=Box]
             {\\}
             {\salien{\texp}{\stau}{\sdelta} \trans \texp}
  %
  \and
  %
  \inferrule*[right=Pair]
             {\source{e_i} \trans \target{e_i}}
             { \spair{\source{e_1}}{\source{e_2}} \trans \tpair{\target{e_1}}{\target{e_2}}}
  %
  \and
  %
  \inferrule*[right=Proj]
             {\sexp \trans \texp}
             {\sproji{\sexp} \trans \tproji{\texp}}
  %
  \and
  %
  \inferrule*[right=Inl]
             {\sexp \trans \texp}
             {\sinl{\sexp} \trans \tinl{\texp}}
  %
  \and
  %
  \inferrule*[right=Inr]
             {\sexp \trans \texp}
             {\sinr{\sexp} \trans \tinr{\texp}}
  %
  \and
  %
  \inferrule*[right=Case]
             {\sexp \trans \texp \\
              \source{e_i}\ssubst{\svar{x}}{\salien{\tvar{x}}{\source{\tau_i}}{\cdot}} \trans \target{e_i}
             }
             {\scasex{\sexp} {\staui}{\sexpi} {\stauii}{\sexpii} \trans \tcasex{\texp} {\texpi} {\texpii}}
  %
  \and
  %
  \inferrule*[right=Abs]
             {\sexp\ssubst{\svar{x}}{\salien{\tvar{x}}{\source{\tau}}{\cdot}} \trans \texp}
             {\sabs{\svar{x}}{\stau}{\sexp} \trans \tabs{\tvar{x}}{\texp}}
  %
  \and
  %
  \inferrule*[right=App]
             {\source{e_i} \trans \target{e_i}}
             {\sapp{\source{e_1}}{\source{e_2}} \trans \tapp{\target{e_1}}{\target{e_2}}}
  %
  \and
  %
  \inferrule*[right=TAbs]
             {\sexp \trans \texp}
             {\stabs{\svar{A}}{\sexp} \trans \texp}
  %
  \and
  %
  \inferrule*[right=TApp]
             {\sexp \trans \texp}
             {\stapp{\sexp}{\stau} \trans \texp}

 %% %
 %%  \and
 %%  %
 %%  \inferrule*[right=Box]
 %%             {\texp \tevals \tval}
 %%             {\salien{\texp}{\stau}{\sdelta} \seval \salien{\tval}{\stau}{\sdelta}}
 %%   %
 %%  \and
 %%  %
 %%  \inferrule*[right=ProjAlien]
 %%             {\\}
 %%             {\salien{\tpair{\texp_{\target{1}}}{\texp_{\target{2}}}}{\sprod{\stau_{\source{1}}}{\stau_{\source{2}}}}{\sdelta}
 %%               \seval
 %%               \salien{\texp_{\target{1}}}{\stau_{\source{1}}}{\sdelta} }
 %%  %
 %%  \and
 %%  %
 %%  \inferrule*[right=CaseLAlien]
 %%             {\\}
 %%             {\scasex{\salien{\tinl{\texp}}{\ssum{\stau_{\source{1}}}{\stau_{\source{2}}}}{\sdelta}}{\stau_{\source{1}}}{\sexp_{\source{1}}}{\stau_{\source{2}}}{\sexp_{\source{2}}}
 %%               \seval
 %%               \sexp_{\source{1}} \ssubst{\svar{x}}{\salien{\texp}{\stau_{\source{1}}}{\sdelta}}}
 %%  %
 %%  \and
 %%  %
 %%  \inferrule*[right=AppAlien]
 %%             {\\}
 %%             {\sapp{\salien{\tabs{\tvar{x}}{\texp}}{\sarrow{\stau_{\source{1}}}{\stau_{\source{2}}}}{\sdelta}}
 %%               {\sexp}
 %%               \seval
 %%               \sabs{\svar{x}}{\sdelta\stau_{\source{1}}}
 %%                    {\kap{\stau_{\source{1}}}{\sdelta}{\svar{x}}
 %%                      {\mfun{\texpp}{\salien{\texp\ssubst{\tvar{x}}{\texpp}}{\stau_{\source{2}}}{\sdelta}}}}
 %%             }
 %%  %
 %%  \and
 %%  %
 %%  \inferrule*[right=TAppAlien]
 %%             {\\}
 %%             {\stapp{\salien{\texp}{\sall{\svar{A}}{\stau}}{\sdelta}}{\staup}
 %%               \seval
 %%               \salien{\texp}{\stau}{\sesubstext{\sdelta}{\svar{A}}{\staup}}
 %%             }
 %%  %
 %%  \and
 %%  %
 %%  \inferrule*[right=AlienOpaque]
 %%             {\stau \not = \sttop}
 %%             {\salien{\topaque{\sexp}}{\stau}{\sdelta}
 %%               \seval
 %%               \sexp
 %%             }
 %%  %
 %%  \and
 %%  %
 %%  \inferrule*[right=AlienCancel]
 %%             {\\}
 %%             {\salien{\talien{\sexp}{\staup}{\sdeltap}}{\stau}{\sdelta}
 %%               \seval
 %%               \sexp}
\end{mathpar}

  \end{framed}
  \caption{Translation}
\end{wrapfigure}

\medskip

\begin{wrapfigure}[htb]{L}{\textwidth}%
    \centering
      \centering
      \begin{framed}
        \[
\begin{array}{r l l}
  \scoerce{\stunit}{\sdelta}{\tunit} & = & \sunit
  \\[0.3em]
  %
  \scoerce{\sarrow{\staui}{\stauii}}{\sdelta}{\tabs{\tvar{x}}{\texp}} & = &
  \sabs{\svar{x}}{\sdelta\staui}{\salien{\texp\tsubst{\tvar{x}}{\prot{\svar{x}}{\staui}{\sdelta}}} {\stauii} {\sdelta}}
  \\[0.3em]
  %
  \scoerce{\sprod{\staui}{\stauii}}{\sdelta}{\tpair{\tvali}{\tvalii}} & = &
  \spair{\salien{\tvali}{\staui}{\sdelta}}{\salien{\tvalii}{\stauii}{\sdelta}}
  \\[0.3em]
  %
  \scoerce{\ssum{\staui}{\stauii}}{\sdelta}{\tinl{\tval}} & = &
  \sinl{\salien{\tval}{\staui}{\sdelta}}
  \\[0.3em]
  %
  \scoerce{\ssum{\staui}{\stauii}}{\sdelta}{\tinr{\tval}} & = &
  \sinr{\salien{\tval}{\stauii}{\sdelta}}
  \\[0.3em]
  % 
  \scoerce{\svar{A}}{\sdelta}{\topaque{\sexp}} & = &
  \sexp \\[0.3em]
  & & \\
  %
  \scoerce{\sall{\svar{A}}{\ssig}}{\sdelta}{\tval} & = &
  \stabs{\svar{A}}{\scoerce{\ssig}{\sdelta}{\tval}}
  \\[0.3em]

  

  %
  %
  %% \scoerce{\sttop}{\sdelta}{\salien{\tval}{\staup}{\sdeltap}} & = &
  %% \tval
  %% \\[0.3em]
\end{array}
\]

      \end{framed}
      \caption{Target to Source Coercion}
      \medskip\medskip\medskip\medskip\medskip
      \centering
      \begin{framed}
        \[
\begin{array}{r l l}
  \tcoerce{\stunit}{\sdelta}{\sunit} & = & \tunit
  \\[0.3em]
  %
  \tcoerce{\sarrow{\staui}{\stauii}}{\sdelta}{\sabs{\svar{x}}{\stau}{\sexp}} & = &
  \tabs{\tvar{x}}{\prot{\sexp\tsubst{\svar{x}}{\salien{\tvar{x}}{\staui}{\sdelta}}} {\stauii} {\sdelta}}
  \\[0.3em]
  %
  \tcoerce{\sprod{\staui}{\stauii}}{\sdelta}{\spair{\sexpi}{\sexpii}} & = &
  \tpair{\prot{\sexpi}{\staui}{\sdelta}}{\prot{\sexpii}{\stauii}{\sdelta}}
  \\[0.3em]
  %
  \tcoerce{\ssum{\staui}{\stauii}}{\sdelta}{\sinl{\sexp}} & = &
  \tinl{\prot{\sexp}{\staui}{\sdelta}}
  \\[0.3em]
  %
  \tcoerce{\ssum{\staui}{\stauii}}{\sdelta}{\sinr{\sexp}} & = &
  \tinr{\prot{\sexp}{\stauii}{\sdelta}}
  \\[0.3em]
  %
  %% \tcoerce{\sall{\svar{A}}{\stau}}{\sdelta}{\stabs{\svar{A}}{\sexp}} & = &
  %% \prot{\sexp\ssubst{\svar{A}}{\sttop}}{\stau\ssubst{\svar{A}}{\sttop} }{\sdelta}
  %% \\[0.3em]
  %% %
  \tcoerce{\svar{A}}{\sdelta}{\sexp} & = & \topaque{\sexp}
  \\[0.3em]
  \tcoerce{\staup}{\sdelta}{\salien{\texp}{\stau}{\sdelta}} & = & \texp
  \hfill \text{if } \stau \equiv \staup
  \\[0.3em]
\end{array}
\]

      \end{framed}
      \caption{Source to Target Coercion}
\end{wrapfigure}

\medskip

\begin{wrapfigure}[htb]{l}{\textwidth}%
  \label{fig:protect}
  %\fbox{$\prot{\sexp}{\stau}{\sdelta}$}
  \begin{framed}
    \[
\prot{\sexp}{\stau}{\sdelta} =
\begin{cases}
  \topaque{\sexp} & \text{if }\stau = \svar{A} \\
  \talien{\sexp}{\stau}{\sdelta} & \text{otherwise}
\end{cases}
\]

  \end{framed}
  \caption{Protect meta-function}
\end{wrapfigure}

\medskip

\begin{wrapfigure}[htb]{l}{\textwidth}%
  \label{fig:typeq}
  %\fbox{$\staui \equiv \stauii$}
  \begin{framed}
    \[
\begin{array}{r l l}
  \stunit \equiv \stunit & \eqdef & \top \\
  %
  \source{A} \equiv \source{B} & \eqdef & \top \\
  %
  \spair{\staui}{\stauii} \equiv
  \spair{\source{\tau_1'}}{\source{\tau_2'}}
  & \eqdef & \staui \equiv \source{\tau_1'} \wedge \stauii \equiv \source{\tau_2'}  \\
  %
  \sarrow{\staui}{\stauii} \equiv \sarrow{\source{\tau_1'}}{\source{\tau_2'}}
  & \eqdef & \staui \equiv \source{\tau_1'} \wedge \stauii \equiv \source{\tau_2'}  \\
  %
  \ssum{\staui}{\stauii} \equiv \ssum{\source{\tau_1'}}{\source{\tau_2'}}
  & \eqdef & \staui \equiv \source{\tau_1'} \wedge \stauii \equiv \source{\tau_2'}  \\
  %
  \stau \equiv \stau' 
  & \eqdef & \bot \hfill \text{otherwise}  \\

\end{array}
\]

  \end{framed}
  \caption{Type equivalence}
\end{wrapfigure}


\medskip

\begin{wrapfigure}[htb]{l}{\textwidth}%
  \label{fig:trans}
  \fbox{$\sexp \trans \texp$} \\
  %$\senv$ is a global parameter 
  \begin{framed}
    \begin{mathpar}
  \inferrule*[right=Unit]
             {\\}
             {\sunit \trans \tunit}
  %
  \and
  \inferrule*[right=Var]
             {\\}
             {\svar{x} \trans \talien{\svar{x}}{\senv\spar{\svar{x}}}{\cdot}}

  %
  \and
  %
  \inferrule*[right=Box]
             {\\}
             {\salien{\texp}{\stau}{\sdelta} \trans \texp}
  %
  \and
  %
  \inferrule*[right=Pair]
             {\source{e_i} \trans \target{e_i}}
             { \spair{\source{e_1}}{\source{e_2}} \trans \tpair{\target{e_1}}{\target{e_2}}}
  %
  \and
  %
  \inferrule*[right=Proj]
             {\sexp \trans \texp}
             {\sproji{\sexp} \trans \tproji{\texp}}
  %
  \and
  %
  \inferrule*[right=Inl]
             {\sexp \trans \texp}
             {\sinl{\sexp} \trans \tinl{\texp}}
  %
  \and
  %
  \inferrule*[right=Inr]
             {\sexp \trans \texp}
             {\sinr{\sexp} \trans \tinr{\texp}}
  %
  \and
  %
  \inferrule*[right=Case]
             {\sexp \trans \texp \\
              \source{e_i}\ssubst{\svar{x}}{\salien{\tvar{x}}{\source{\tau_i}}{\cdot}} \trans \target{e_i}
             }
             {\scasex{\sexp} {\staui}{\sexpi} {\stauii}{\sexpii} \trans \tcasex{\texp} {\texpi} {\texpii}}
  %
  \and
  %
  \inferrule*[right=Abs]
             {\sexp\ssubst{\svar{x}}{\salien{\tvar{x}}{\source{\tau}}{\cdot}} \trans \texp}
             {\sabs{\svar{x}}{\stau}{\sexp} \trans \tabs{\tvar{x}}{\texp}}
  %
  \and
  %
  \inferrule*[right=App]
             {\source{e_i} \trans \target{e_i}}
             {\sapp{\source{e_1}}{\source{e_2}} \trans \tapp{\target{e_1}}{\target{e_2}}}
  %
  \and
  %
  \inferrule*[right=TAbs]
             {\sexp \trans \texp}
             {\stabs{\svar{A}}{\sexp} \trans \texp}
  %
  \and
  %
  \inferrule*[right=TApp]
             {\sexp \trans \texp}
             {\stapp{\sexp}{\stau} \trans \texp}

 %% %
 %%  \and
 %%  %
 %%  \inferrule*[right=Box]
 %%             {\texp \tevals \tval}
 %%             {\salien{\texp}{\stau}{\sdelta} \seval \salien{\tval}{\stau}{\sdelta}}
 %%   %
 %%  \and
 %%  %
 %%  \inferrule*[right=ProjAlien]
 %%             {\\}
 %%             {\salien{\tpair{\texp_{\target{1}}}{\texp_{\target{2}}}}{\sprod{\stau_{\source{1}}}{\stau_{\source{2}}}}{\sdelta}
 %%               \seval
 %%               \salien{\texp_{\target{1}}}{\stau_{\source{1}}}{\sdelta} }
 %%  %
 %%  \and
 %%  %
 %%  \inferrule*[right=CaseLAlien]
 %%             {\\}
 %%             {\scasex{\salien{\tinl{\texp}}{\ssum{\stau_{\source{1}}}{\stau_{\source{2}}}}{\sdelta}}{\stau_{\source{1}}}{\sexp_{\source{1}}}{\stau_{\source{2}}}{\sexp_{\source{2}}}
 %%               \seval
 %%               \sexp_{\source{1}} \ssubst{\svar{x}}{\salien{\texp}{\stau_{\source{1}}}{\sdelta}}}
 %%  %
 %%  \and
 %%  %
 %%  \inferrule*[right=AppAlien]
 %%             {\\}
 %%             {\sapp{\salien{\tabs{\tvar{x}}{\texp}}{\sarrow{\stau_{\source{1}}}{\stau_{\source{2}}}}{\sdelta}}
 %%               {\sexp}
 %%               \seval
 %%               \sabs{\svar{x}}{\sdelta\stau_{\source{1}}}
 %%                    {\kap{\stau_{\source{1}}}{\sdelta}{\svar{x}}
 %%                      {\mfun{\texpp}{\salien{\texp\ssubst{\tvar{x}}{\texpp}}{\stau_{\source{2}}}{\sdelta}}}}
 %%             }
 %%  %
 %%  \and
 %%  %
 %%  \inferrule*[right=TAppAlien]
 %%             {\\}
 %%             {\stapp{\salien{\texp}{\sall{\svar{A}}{\stau}}{\sdelta}}{\staup}
 %%               \seval
 %%               \salien{\texp}{\stau}{\sesubstext{\sdelta}{\svar{A}}{\staup}}
 %%             }
 %%  %
 %%  \and
 %%  %
 %%  \inferrule*[right=AlienOpaque]
 %%             {\stau \not = \sttop}
 %%             {\salien{\topaque{\sexp}}{\stau}{\sdelta}
 %%               \seval
 %%               \sexp
 %%             }
 %%  %
 %%  \and
 %%  %
 %%  \inferrule*[right=AlienCancel]
 %%             {\\}
 %%             {\salien{\talien{\sexp}{\staup}{\sdeltap}}{\stau}{\sdelta}
 %%               \seval
 %%               \sexp}
\end{mathpar}

  \end{framed}
  \caption{Translation}
  \label{fig:translationlong}
\end{wrapfigure}
